% Abstract page — uses the \abstract{} and \keywords{} commands from unicam_thesis.cls

\abstract{%
Renal cell carcinoma (RCC) accounts for approximately 400,000 new diagnoses annually worldwide, with metastatic disease carrying a five-year survival rate of roughly 15\%. Understanding why individual patients respond differently to immunotherapy and targeted therapy remains a key challenge in clinical oncology.

This work presents an agent-based model (ABM) of the RCC tumor microenvironment comprising 18 distinct agent types---including tumor cells, six T~cell subtypes, macrophages, dendritic cells, NK cells, blood vessels, and sex hormone agents---interacting on a three-dimensional discrete grid. Each tumor cell carries a 16-gene DNA profile that governs proliferation, immune evasion, and treatment sensitivity through codon-level mutation mechanics. A continuous glucose concentration field, governed by Fick's diffusion law, overlays the agent grid, coupling metabolic competition (the Warburg effect) to both tumor growth and immune effector function.

The model was ported from Mesa to Repast4Py to enable MPI-based parallel execution and integrates PD-1/PD-L1 checkpoint interactions, ICI and TKI treatment mechanisms, and sex hormone modulation of immune responses across 17 effect channels. The 87 tunable weight parameters are calibrated against ARON-1 clinical trial data using Bayesian optimization with Optuna.

A 160-simulation multi-seed validation study reveals sex-stratified treatment responses: female virtual patients receiving combination ICI+TKI therapy achieve the highest observed survival rate (65\%), consistent with emerging clinical evidence for sex-dependent treatment efficacy. Statistical analysis using Mann-Whitney U tests characterizes the stochastic variability inherent in agent-based tumor--immune dynamics.

These results demonstrate that mechanistic agent-based models incorporating metabolic, genetic, and hormonal subsystems can generate clinically relevant predictions when calibrated against trial data, contributing to ongoing efforts in computational precision oncology.%
}

\keywords{%
Agent-based modeling, renal cell carcinoma, tumor microenvironment, glucose metabolism, Warburg effect, immune checkpoint inhibitors, tyrosine kinase inhibitors, sex hormones, Repast4Py, Bayesian optimization, ARON-1%
}
